\chapter{Intégrale et théorème fondamental de l’analyse}
\section*{Exercice 19.1 -- Aire signée}
[Base] Calculer $\int_{-1}^1x\,dx$ puis l'aire entre sa courbe et l'axe sur ce segment.

\section*{Exercice 19.2 -- Borne variable}
[Méthode] Dériver $F(x)=\int_1^{x^2}e^{-t^2}\,dt$.

\section*{Exercice 19.3 -- Valeur moyenne}
[Application] Calculer la valeur moyenne de $t^2$ sur $[0,3]$ et trouver un point où elle est atteinte.

\section*{Exercice 19.4 -- Positivité stricte}
[Preuve] Si $f$ est continue, positive ou nulle sur $[a,b]$, $a<b$, et $\int_a^bf=0$, montrer que $f$ est identiquement nulle.

\section*{Exercice 19.5 -- Orientation}
[Base] Sachant $\int_0^2f=5$ et $\int_0^5f=1$, calculer $\int_2^5f$ et $\int_5^2f$.

\section*{Exercice 19.6 -- Comparaison}
[Raisonnement] Pour $0\le x\le1$, encadrer $\int_0^1e^{-x^2}\,dx$ avec des constantes simples, puis améliorer la borne inférieure par $e^{-u}\ge1-u$.
